% SPDX-FileCopyrightText: Copyright © 2026 Idiap Research Institute <contact@idiap.ch>
% SPDX-FileContributor: Andre Freitas <andre.freitas@idiap.ch>
% SPDX-License-Identifier: MIT
% Neuro-symbolic AI Group
% File Documentation
%
% Motivation:
%   Protect Mathformer behavior while the framework evolves across Lean, ATP, KB, SymPy, docs, and skills.
%
% Intent:
%   Encode regression checks that catch parser breakage, semantic drift, and missing workflow evidence.
%
% Function:
%   Contains executable tests, fixtures, and manifests used for local and CI validation.
%
% Maintenance Notes:
%   Keep this header synchronized with major responsibility changes, and prefer updating tests or examples when behavior changes.

% mathformer-case-id: wiki_contrapositive
% expected-lean-name: WikiLevelGold.wiki_contrapositive
% controlled-suite: wiki_level_latex_to_lean
% difficulty: very_easy
\documentclass[11pt]{article}
\usepackage[T1]{fontenc}
\usepackage[utf8]{inputenc}
\usepackage{lmodern}
\usepackage{amsmath,amssymb,amsthm}

\newtheorem{theorem}{Theorem}

\title{The Contrapositive}
\author{}
\date{}

\begin{document}

\maketitle

\section*{Controlled Origin Annotation}
This is a synthetic controlled case written in a Wikipedia-style exposition.
It is intended to test whether the LaTeX-to-Lean pipeline preserves hypotheses,
local bindings, negation, and implication order.

\section*{Statement and Proof}

\begin{theorem}[Contrapositive]
Let $P$ and $Q$ be propositions. If $P$ implies $Q$, then not $Q$ implies
not $P$.
\end{theorem}

\begin{proof}
Assume first that $h : P \to Q$. To prove the desired implication, assume
$\neg Q$ and suppose, for contradiction, that $P$ holds. Applying $h$ to the
proof of $P$ gives a proof of $Q$, contradicting $\neg Q$. Therefore $P$ is
impossible, so $\neg P$ holds. Hence $\neg Q \to \neg P$.
\end{proof}

\section*{Lean Target}
Formalize the theorem for arbitrary propositions $P$ and $Q$ as
\[
  (P \to Q) \to (\neg Q \to \neg P).
\]
The expected Lean theorem is \texttt{WikiLevelGold.wiki\_contrapositive}.

\end{document}
